% Abstract page, following the Word paragraph settings
{\centering\textbf{ABSTRACT}\par}
\vspace{34.5pt}
{\centering\fontsize{14}{24.15}\selectfont\bfseries \thesistitleA\\ \thesistitleB\par}
\vspace{34.5pt}

This thesis examines government support for crop disaster insurance in South Korea. I assess premium subsidies, administrative cost subsidies, and state reinsurance, and compare agricultural outcomes around insurance introduction. Each instrument raises a different evaluation question: who can obtain protection and who receives support, what delivery costs achieve, and how losses are shared between private insurers and the state. I then compare six annual crops with crops not yet exposed to insurance and compare rice planted area in early and later pilot municipalities. The mean production contrast across the six crops is −0.002 log points, with an exploratory 95 percent interval of [−0.318, 0.315]. The estimates vary with reference periods, comparison crops, and weights, and area and production show differences in pre-pilot growth. The rice estimate is −0.021 log points, with a 95 percent interval of [−0.064, 0.022], and changes with sample composition. These comparisons leave substantial uncertainty about production responses. Together, the institutional and empirical analyses show why the assessment of public support requires evidence on access, distribution, service quality, and fiscal exposure as well as production.

\vspace{20.7pt}
{\wordonehalf\noindent\emph{Keywords}: crop disaster insurance; government intervention; premium subsidies; state reinsurance; South Korea\par}
\clearpage
